% Set from report/draft.md section 4. The prose is the maintainer's; this file is
% its transit into LaTeX and nothing in it was rewritten.
\section{Findings}

\subsection{Summary}

Thirteen experiments produced a set of results that read at first as unrelated:
some about the synthetic stimulus, some about the gaze policy, one about closing
the loop. They are better understood as one claim measured at three levels.

\begin{quote}
\textbf{Availability dominates selection.} What can be measured determines the
result. How cleverly one chooses among what can be measured is a second-order
refinement. This holds whether the thing limiting availability is the stimulus,
the sensor's field, or the geometry of occlusion.
\end{quote}

Each experiment was designed to answer the question in front of it --- does the
objective matter, does variance help, does verging acquire anything --- and the
pattern across them became legible only once the instrument was characterised
well enough that the answers could be trusted. \S4.2 asks what selection buys;
\S4.3 asks what changes availability.

Two consequences of \S3.3 have to be carried forward, because both were
reversals.

Conclusions already drawn had to be re-read, because results measured on the
fixture were pooled over a pixel population dominated by the artefact.

And the obvious correction was also wrong. Restricting the analysis to pixels
far from any discontinuity made every policy difference vanish, which looked
like a clean answer: the effects had been artefactual. On rendered data they
returned, stronger. Removing discontinuities does not isolate \emph{clean}
measurement --- it isolates \emph{easy} measurement, and on smooth textured
surface there is nothing to allocate. The band where the artefact lived is also
the band where the question lives, which is why the artefact was so damaging and
why removing it by masking answers nothing.

A stimulus, once characterised, also has an operating range. Sweeping the
fraction of occluded surface showed the allocation question is sharpest near a
tenth of the surface occluded --- and, independently, that the acquisition gain
measured later peaks at the same level. Above roughly a sixth the stimulus stops
discriminating between policies altogether, and not merely because the spread
grows: the arms' error means themselves converge, from a spread of 0.51 at a
tenth to 0.02 at a sixth, as the scene's own difficulty saturates the estimator.
That range then sized every experiment after it.

A stimulus is an instrument. Four of thirteen experiments went to establishing
that, and none of the results after them would have been readable otherwise.

\subsection{What selection buys}

The gaze policy proposes a next fixation by maximising something over the
current belief. The natural candidate is posterior variance: look where you are
least certain.

\textbf{The choice of objective does not matter.} Two objectives were compared
--- variance, and the expected reduction in variance a fixation would actually
produce, integrated over the field. They never once selected the same pixel: an
agreement rate of exactly zero across 144 comparisons, and zero again across a
further 144 for a second pairing. Their picks sit a median 115 pixels apart,
more than three fovea widths --- different regions of the image, not
neighbouring cells. And their results were indistinguishable under the
experiment's pre-registered bar, with a sign test agreeing there is no effect.

The fields are moderately rank-correlated, so this is a statement about the
argmax rather than about the objectives being unrelated. It is also a statement
about one fixture, with both objectives evaluated at the same state on the same
candidate grid, so that only the objective varies.

\textbf{What matters is not looking twice.} Posterior variance conflates two
situations: \emph{nothing has been measured here yet}, and \emph{measurement
here is known to be uninformative}. The second is a fixed point --- looking
again does not reduce it. Without a mechanism to suppress recently-visited
locations, the policy locks: in one representative run it visited eight distinct
locations and then spent its remaining ten fixations on a single pixel, where
the matcher had confidently hallucinated a disparity and the confidence
machinery had correctly refused it.

\begin{figure}[htbp]
  \centering
  \includegraphics[width=\textwidth]{scanpath.pdf}
  \caption{Scanpath of the ported baseline under argmax-posterior-variance
  selection with no inhibition of return, on the analytic fixture at seed 0.
  Eighteen fixations, nine of them distinct: the loop visits eight pixels once
  each, reaches pixel $(170, 94)$ at step 8, and spends its remaining ten
  fixations there. Step 6 lands at $(169, 94)$, one pixel above the lockup
  pixel, so its marker sits beneath the highlighted one at this scale.
  \emph{Measured}, and regenerated for this figure from committed code at a
  pinned seed: \texttt{bio-001} and \texttt{bio-002} in
  \texttt{docs/inherited-measurements.yaml}, pinned by
  \texttt{tests/test\_loop\_port.py}. This is the predecessor's loop, ported
  bit-identically (\texttt{bio-009}); it is inherited behaviour and a baseline,
  not a result about anything new.}
  \label{fig:scanpath}
\end{figure}

\textbf{Knowing where measurement is possible beats knowing how well.} A policy
informed only by the validity mask --- no variance at all, just \emph{don't look
where nothing can be measured} --- captures the overwhelming majority of the
available benefit. On median error, on a stimulus with effectively no occlusion,
knowing \emph{where} the data is was worth 65\texttimes{} what knowing \emph{how
uncertain} it is added on top. On the tail the ratio narrows to 3.5:1, and it
falls as occlusion rises. Variance is a real refinement, and a second-order one.

That refinement was nearly lost to the instrument. On the fixture, variance
appeared to contribute nothing once the artefact band was excluded. On rendered
data it recovered, several times stronger than the fixture had ever shown. The
fixture was understating the saliency channel, not inflating it --- the opposite
of the concern that prompted the check.

\textbf{Foveal weighting, on a uniformly sampled sensor, is strictly lossy.}
Given a complete spherical capture --- every direction present, nothing outside
the field --- a gaze-dependent foveal weight lost to uniform weighting on both
metrics, in every band grouping, on every one of eight seeds. The mechanism is
measured: the solid angle receiving more than one percent of peak weight is the
whole sphere for uniform weighting and about a fifth of it for foveal, so the
foveal arm leaves four fifths of the sphere at the prior. The reason is
structural rather than empirical --- a foveal weight models an \emph{acuity
gradient}, and a uniformly sampled sensor does not have one. Weighting by gaze
can only discard information already paid for. It cannot buy anything back.

This does not mean allocation is worthless. It means allocation implemented as
weighting, on a sensor with no acuity gradient, has nothing to work with.

\subsection{What changes availability}

If selection buys little, the remaining hypothesis is that active vision works
by changing what is available rather than by choosing among it.

\textbf{Deliberately verging to acquire made things worse.} Narrowing the
disparity search around the current vergence estimate, so that re-verging brings
new depths into range, cost 7.5\texttimes{} in median accuracy. Verging to
re-acquire did not recover it: it halved the tail and bought 21.7\% more
coverage, at 2.6\texttimes{} the hypothesis count and nineteen front-end calls
against one.

The mechanism is the interesting part, and it is not specific to vergence. The
estimator computes a statistic over \emph{valid} measurements in a foveal
window, and a target outside the search range either fails validity or clamps to
the edge --- so the pixels carrying the error are the ones the validity mask has
already deleted. At one measured fixation the target sat at nineteen pixels of
disparity while the estimator's window reported eight, and the controller verged
away from it, from ten pixels to six over three fixations. Across the run it
partially recovers, and it is better described as badly steered than as stuck.
\S5 returns to this as the finding with the widest reach beyond stereo.

\textbf{Closing the perception-action loop does help.} Re-rendering at each new
fixation, re-matching, and accumulating into a head-centred belief beats holding
a single capture and re-weighting it, at 41\texttimes{} the cost in wall-clock
seconds --- of which the renders are 91\%.

It helped at every occlusion level tested, consistently in direction: the moving
arm was better in all sixteen band-by-level cells. It was not consistent in
verdict. Per band the results are level-dependent and non-monotone, and the
effect at depth discontinuities crosses the bar at the lowest and highest
occlusion levels but not at the two between. A consolidation run at three times
the seeds also revised the original margins downward, because the smaller sample
had underestimated the spread. The conclusion survived; its margin did not.

Decomposing the gain is where the result becomes specific. It splits into an
\emph{accuracy} term, on the cells both arms measured, and a \emph{coverage}
term, on the cells only the moving sensor ever reached. Coverage exceeded
accuracy by at least 33\texttimes{} at every level and by 341\texttimes{} at the
highest, in a decomposition that is exact --- the two terms sum to the total at
every level. The closed loop's value is not that it estimates better; it is that
it sees the scene, and the fixed sensor does not.

\begin{figure}[htbp]
  \centering
  \includegraphics[width=\textwidth]{coverage.pdf}
  \caption{The two arms of exp010/exp011 over four occlusion levels, 24 seeds
  each. \emph{Left}: the fraction of the belief each arm ever measured.
  \emph{Right}: exp011's question-3 decomposition of the all-cells gain into an
  accuracy term, on the cells both arms reached, and a coverage term, on the
  cells only the moving arm reached; note the logarithmic axis. \emph{Measured}:
  exp011, on the pre-registered yardstick of mean absolute error over all
  76\,800 cells with unmeasured cells left at the prior
  (\texttt{experiments/exp011\_consolidation/results.json}, 192 rows;
  \texttt{bio-076}). The decomposition is exact --- the two terms sum to the
  total at every level --- and this figure rebuilds it from the per-seed rows
  and asserts the result against the values exp011 published before drawing.}
  \label{fig:coverage}
\end{figure}

The accuracy term was then narrowed by an experiment whose premise turned out to
be false. Both arms re-estimate vergence at every fixation and both apply the
same foveal weight at the same locations, so they differ in exactly one respect:
whether the images are fresh. Away from depth discontinuities, freezing the
linearisation pedestal costs nothing --- so the advantage there comes from the
fresh captures rather than from re-linearising about each new fixation. At
discontinuities the experiment cannot separate the two, and one confound remains
unexcluded: the stationary arm's measurements are far more correlated with each
other than the moving arm's. The attribution to re-acquisition is established on
smooth surface and open where the geometry is hard.

That is the project's positive result, and it is a narrow one. A second
viewpoint helps. Choosing it cleverly does not.

\textbf{And on a complete spherical sensor, acquisition disappears entirely.}
The eyes rotate about centres fixed in the head; a rotation about a fixed
optical centre does not change what a complete sphere sees. Re-rendering at a
new fixation returns an identical capture. There is nothing to acquire.

\begin{figure}[htbp]
  \centering
  \includegraphics[width=\textwidth]{rig.pdf}
  \caption{Why a saccade acquires nothing on a complete spherical sensor.
  \emph{Top}: two eye centres fixed in the head at $\pm b/2$, each carrying a
  sensor that has already sampled every direction; the gaze rotates and the
  capture is unchanged, because the rotation moves the \emph{frame} and not the
  \emph{data}. \emph{Bottom}, for contrast: on a pinhole sensor the field of
  view rotates with the gaze, and the crescent it uncovers is what acquisition
  means. A conceptual diagram --- it draws a geometric argument, so every
  factual mark traces to a ledger entry, and the script asserts those entries
  exist before drawing. \texttt{fc-013}: the centres are fixed and
  \texttt{eye\_rotations} returns rotations, not translations, so ``a complete
  capture has no hemisphere to hide''. \texttt{bio-081}: a sphere loses no grid
  at any saccade amplitude. \texttt{bio-082}: two orientations, identity and
  yawed $37^\circ$/pitched $19^\circ$, reproduce the same analytic function of
  direction to 80\,\textmu m, measured against a real Blender. \textbf{The
  qualifier is load-bearing}: this holds for a sensor sampled \emph{uniformly}.
  Under variable resolution the fovea resolves what the periphery only sampled
  coarsely, and refixation buys something again (\texttt{od-004}).}
  \label{fig:rig}
\end{figure}

\begin{figure}[htbp]
  \centering
  \includegraphics[width=\textwidth]{spherical.pdf}
  \caption{Spherical stereo geometry. \emph{Left}: with the baseline along
  $+x$ the two epipoles sit at $\pm x$, so every epipolar great circle passes
  through them and the family of such circles is exactly the meridians about the
  baseline --- which is why rows of the equirectangular sensor are epipolar and
  the matcher's horizontal shift is a shift in colatitude $\theta$.
  \emph{Right}: the sine rule $r_L = b\,\sin\theta_R/\sin d$ on the triangle
  (left centre, right centre, point) makes the angular disparity depend on
  $\theta$, the angle from the baseline, in a way the pinhole relation $f b / z$
  has no analogue for. \emph{Derived}, closed-form, from the geometry documented
  at \texttt{src/bio3dvision/sampling.py:242--264}, whose stated consequence ---
  at 3\,m the disparity is 5.8\texttimes{} larger at $\theta = 90^\circ$ than at
  $10^\circ$, here 5.76\texttimes{} --- this figure asserts before drawing. The
  relation was checked against a render to $10^{-14}$ relative
  (\texttt{bio-077}).}
  \label{fig:spherical}
\end{figure}

That is not a limitation of the implementation. It is the geometry, and it sets
up the question \S5 takes as its subject.
